La Terre reçoit du Soleil une intensité de \qty{\v{S}}{W/m^2} (constante solaire) et réfléchit une fraction \num{\v{al}} de ce rayonnement (albédo). On suppose que la Terre rayonne comme un corps noir, uniformément sur toute sa surface, et on néglige l'atmosphère.

\begin{enumerate}[label=\alph*)]
    \item Expliquez pourquoi la puissance absorbée par mètre carré de surface terrestre, en moyenne, vaut $S(1 - a)/4$, puis calculez-la.
    \item Déterminez la température d'équilibre de la Terre.
    \item La température moyenne réelle de la surface est d'environ \qty{288}{K}. Quelle est la cause de l'écart ?
\end{enumerate}
